\chapter*{Preface}
\addcontentsline{toc}{chapter}{Preface}

A service answers five requests in 112, 98, 105, 240, and 101 ms. The mean is
131.2 ms; the median (the middle value once sorted) is 105 ms. Without the
240, the other four have mean 104 ms and median 103 ms. One slow request
pulled the mean up by 27 ms and moved the median by only 2 ms
(Figure~\ref{fig:pre-mean-median}).

\begin{figure}[H]
\centering
\begin{tikzpicture}
\begin{axis}[name=main, width=8.4cm, height=4.8cm,
  xmin=94, xmax=136, ymin=0.2, ymax=3.4,
  axis lines*=left, axis y line=none,
  xtick={95,100,...,135}, xlabel={latency (ms)},
  tick label style={font=\scriptsize}, label style={font=\small},
  clip=false]
\node[anchor=east, font=\footnotesize, align=right] at (axis cs:93,2.2)
  {all five\\ requests};
\node[anchor=east, font=\footnotesize, align=right] at (axis cs:93,0.9)
  {without the\\ 240 ms one};
\draw[targetgreen, very thick] (axis cs:105,1.95) -- (axis cs:105,2.45);
\draw[densityred, very thick] (axis cs:131.2,1.95) -- (axis cs:131.2,2.45);
\draw[targetgreen, very thick] (axis cs:103,0.65) -- (axis cs:103,1.15);
\draw[densityred, very thick] (axis cs:104,0.65) -- (axis cs:104,1.15);
\addplot[only marks, mark=*, mark size=2.2pt, color=distblue]
  coordinates {(98,2.2) (101,2.2) (105,2.2) (112,2.2)
               (98,0.9) (101,0.9) (105,0.9) (112,0.9)};
\node[anchor=north east, font=\footnotesize, targetgreen, inner xsep=1pt]
  at (axis cs:105.6,1.9) {median 105};
\node[anchor=north east, font=\footnotesize, densityred, inner xsep=1pt]
  at (axis cs:131.8,1.9) {mean 131.2};
\node[anchor=west, font=\footnotesize, align=left] at (axis cs:114,0.9)
  {{\color{targetgreen}median 103}, {\color{densityred}mean 104}:\\ almost equal};
\draw[densityred, -{Stealth[length=2mm]}, thick]
  (axis cs:104,2.75) -- (axis cs:131.2,2.75);
\node[anchor=south, font=\footnotesize, densityred] at (axis cs:117.6,2.8)
  {one slow request drags the mean up 27 ms};
\end{axis}
\begin{axis}[at={(main.south east)}, anchor=south west, xshift=0.7cm,
  width=2.4cm, height=4.8cm,
  xmin=235, xmax=245, ymin=0.2, ymax=3.4,
  axis lines*=left, axis y line=none,
  xtick={240}, tick label style={font=\scriptsize}, clip=false]
\addplot[only marks, mark=*, mark size=2.2pt, color=trapred] coordinates {(240,2.2)};
\node[anchor=south, font=\footnotesize, trapred, align=center] at (axis cs:240,2.35)
  {slow request\\ 240 ms};
\end{axis}
\node[font=\footnotesize, gray] at ($(main.south east)+(0.35cm,0.0cm)$) {//};
\end{tikzpicture}
\caption{The five latencies from the opening example, one dot per request.
Compare the two rows: removing the single 240 ms request moves the red mean
line by 27 ms but the green median line by only 2 ms, because the median looks
only at the middle value. The axis is broken between 136 and 235 ms, where no
request falls, so the four fast dots stay far enough apart to count.}
\label{fig:pre-mean-median}
\end{figure}

Which number should go on the dashboard? If
tomorrow's mean is 125 ms, did the service get faster, or did it just draw a
different five requests? How many requests would you need before a 5 ms change
is real rather than luck? This book exists to answer questions of that kind,
with numbers you can compute and reasons you can check.

\section*{Who this book is for}

The reader we have in mind is an engineer, analyst, or student who uses
statistics every week and wants to stop taking it on faith. You can read code
and algebra. You may have computed a standard deviation or run a $t$-test, but
you could not say why the standard deviation divides by $n-1$, or why averaging
400 measurements is only twice as precise as averaging 100.

The level is that of a solid first course: the content of Khan Academy's
statistics track and AP Statistics, plus the statistics that machine learning
(ML) texts assume you already know: \emph{An Introduction to Statistical
Learning} (ISLR), or the opening chapters of Bishop's \emph{Pattern
Recognition and Machine Learning} and Goodfellow, Bengio, and Courville's
\emph{Deep Learning}. Nothing beyond algebra and a little calculus (sums, a derivative, the
idea of an integral as area) is needed.

\section*{What you will be able to do}

By the end you should be able to:
\begin{itemize}
  \item read a latency dashboard and say what p50, p95, and p99 mean (the
    values that 50\%, 95\%, and 99\% of requests stay under), and when the
    mean misleads;
  \item explain why the sample standard deviation (the typical distance of a
    value from the mean) divides by $n-1$, and why the standard error of a
    mean (the typical distance of the mean itself from the truth) is
    $\sigma/\sqrt{n}$;
  \item build and interpret confidence intervals for a mean, a proportion, and
    a median, including a bootstrap interval for any statistic;
  \item run an A/B test, read its $p$-value correctly (how often data at
    least this extreme would turn up if there were no real difference), and
    choose its sample size in advance;
  \item recognise multiple-testing, peeking, and confounding before they
    produce a false claim;
  \item fit a linear or logistic regression, read every number in its output,
    and check its residuals;
  \item read notation such as $\E_{x\sim p}[\ell(\theta;x)]$ in a machine
    learning paper (the average loss $\ell$ of parameters $\theta$ over data
    $x$ drawn from a distribution $p$) and know which statistical idea it
    encodes.
\end{itemize}

\section*{The road map}

The book has sixteen chapters in five parts. Each part uses the tools of the
one before it. Figure~\ref{fig:pre-roadmap} shows the path.

\begin{figure}[H]
\centering
\begin{tikzpicture}[
  font=\scriptsize,
  stage/.style={draw, rounded corners, fill=#1!18, minimum width=2.3cm, inner xsep=2pt,
    minimum height=1.2cm, align=center, font=\scriptsize\bfseries},
  ch/.style={draw=#1!70!black, fill=white, minimum width=2.35cm, inner xsep=2pt,
    minimum height=1.2cm, align=center, rounded corners=2pt},
  flow/.style={-{Stealth[length=2.2mm]}, thick, gray!70!black}
]
\node[stage=sourceblue] (s1) at (0,0) {I.\\Describe};
\node[ch=sourceblue, right=0.2cm of s1] (c1) {1 Data and\\distributions};
\node[ch=sourceblue, right=0.12cm of c1] (c2) {2 Center\\and spread};
\node[ch=sourceblue, right=0.12cm of c2] (c3) {3 Quantiles,\\box plots,\\shape};
\node[ch=sourceblue, right=0.12cm of c3] (c4) {4 Two\\variables,\\correlation};
\node[stage=targetgreen] (s2) at (0,-1.6) {II. Sample\\$\to$ population};
\node[ch=targetgreen, right=0.2cm of s2] (c5) {5 Sampling and\\study design};
\node[ch=targetgreen, right=0.12cm of c5] (c6) {6 Sampling\\distributions,\\SE};
\node[ch=targetgreen, right=0.12cm of c6] (c7) {7 Confidence\\intervals};
\node[ch=targetgreen, right=0.12cm of c7] (c8) {8 The\\bootstrap};
\node[stage=quotientorange] (s3) at (0,-3.2) {III. Test\\and compare};
\node[ch=quotientorange, right=0.2cm of s3] (c9) {9 Hypothesis\\testing};
\node[ch=quotientorange, right=0.12cm of c9] (c10) {10 Comparing\\groups and\\categories};
\node[ch=quotientorange, right=0.12cm of c10] (c11) {11 Effect sizes,\\multiple\\testing};
\node[stage=structurepurple] (s4) at (0,-4.8) {IV.\\Models};
\node[ch=structurepurple, right=0.2cm of s4] (c12) {12 Simple\\linear\\regression};
\node[ch=structurepurple, right=0.12cm of c12] (c13) {13 Multiple and\\logistic\\regression};
\node[ch=structurepurple, right=0.12cm of c13] (c14) {14 Likelihood,\\estimation,\\Bayes};
\node[ch=structurepurple, right=0.12cm of c14] (c15) {15 Statistics\\for machine\\learning};
\node[stage=gray] (s5) at (0,-6.4) {V.\\Reference};
\node[draw, dashed, rounded corners, fill=gray!8, minimum width=9.76cm, inner xsep=2pt,
  minimum height=0.8cm, align=center, right=0.2cm of s5] (c16)
  {16 Appendix: formula sheet and decision guide (use any time)};
\foreach \a/\b in {s1/s2, s2/s3, s3/s4, s4/s5} \draw[flow] (\a) -- (\b);
\foreach \a/\b in {c1/c2, c2/c3, c3/c4, c5/c6, c6/c7, c7/c8, c9/c10, c10/c11,
  c12/c13, c13/c14, c14/c15} \draw[flow] (\a) -- (\b);
\end{tikzpicture}
\caption{The sixteen chapters in five parts. Read top to bottom: you describe
data, then learn how a sample speaks for a population, then use that to test
claims, then build models and carry them into machine learning.
Chapter 16 is a reference you can consult at any point.}
\label{fig:pre-roadmap}
\end{figure}

\textbf{Part I (Chapters 1--4)} works with data you already have: histograms,
means, medians, standard deviations, quantiles, and correlation. No probability
is needed yet.

\textbf{Part II (Chapters 5--8)} asks what a sample says about the population
it came from. This is where the standard error and the confidence interval
appear, and where the bootstrap (resampling your own data many times to see how much
a statistic wobbles) lets a computer do the algebra for you.

\textbf{Part III (Chapters 9--11)} turns intervals into decisions: $p$-values,
$t$-tests, A/B tests, chi-square tests, effect sizes, and the ways testing
goes wrong.

\textbf{Part IV (Chapters 12--15)} fits models. Regression first, then the
general principles behind it (likelihood, Bayes), then the statistics inside
machine learning: train/test splits, cross-validation, bias and variance, and
loss functions as expectations.

\textbf{Part V (Chapter 16)} is a formula sheet and a decision guide for
choosing a method. Keep it open while you read.

\section*{The five ideas that run the whole book}

Almost every chapter is an application of five ideas.
Figure~\ref{fig:pre-five-ideas} draws them.

\begin{enumerate}
  \item \textbf{Variability.} Measurements differ. The standard deviation
    measures by how much (Chapter 2).
  \item \textbf{Sampling distribution.} A statistic computed from a sample,
    such as the sample mean $\xbar$, would come out differently on another
    sample. Its distribution over all possible samples is the object
    inference is about (Chapter 6).
  \item \textbf{$\SE \propto 1/\sqrt{n}$.} If single values have standard
    deviation $\sigma$, the spread of $\xbar$ over repeated samples of size
    $n$ is $\sigma/\sqrt{n}$, called the standard error ($\SE$). To halve the
    uncertainty you need four times the data (Chapter 6).
  \item \textbf{Uncertainty intervals.} Report an estimate as
    estimate $\pm$ (about 2) $\times \SE$. The ``95\%'' is a promise about
    the procedure: over many repeated samples, about 95\% of the intervals
    built this way contain the true value (Chapters 7 and 8).
  \item \textbf{Models are signal plus noise.} $y = f(x) + \varepsilon$: the
    outcome $y$ is a systematic part $f(x)$ plus random noise $\varepsilon$.
    Regression, tests, and ML all estimate $f$ while measuring how much the
    noise could have misled them (Chapters 12--15).
\end{enumerate}

\begin{figure}[H]
\centering
\begin{tikzpicture}[font=\scriptsize,
  panel/.style={draw=gray!60, rounded corners, minimum width=3.95cm,
    minimum height=3.5cm}]
\begin{scope}[shift={(0,0)}]
\node[panel] at (1.9,1.0) {};
\node[font=\scriptsize\bfseries] at (1.9,2.5) {1. Variability};
\draw[->] (0.2,0.3) -- (3.6,0.3);
\foreach \x/\h in {0.6/1,1.0/1,1.3/1,1.3/2,1.7/1,1.7/2,1.7/3,2.0/1,2.0/2,
  2.3/1,2.3/2,2.7/1,3.2/1}
  \fill[sourceblue] (\x,0.3+0.22*\h) circle (0.08);
\draw[gray, dashed] (1.83,0.3) -- (1.83,1.35);
\draw[densityred, thick, |-|] (1.13,1.35) -- (2.53,1.35);
\node[densityred] at (1.83,1.62) {mean $\pm$ 1 SD};
\node[align=center] at (1.9,-0.15) {13 values that differ;\\ the SD says by how much};
\end{scope}
\begin{scope}[shift={(4.2,0)}]
\node[panel] at (1.9,1.0) {};
\node[font=\scriptsize\bfseries] at (1.9,2.5) {2. Sampling distribution};
\draw[->] (0.2,0.3) -- (3.6,0.3) node[above left] {$\xbar$};
\foreach \x/\h in {0.8/0.15,1.1/0.4,1.4/0.85,1.7/1.3,2.0/1.3,2.3/0.85,
  2.6/0.4,2.9/0.15}
  \fill[histfill, draw=distblue] (\x,0.3) rectangle (\x+0.3,0.3+\h);
\draw[densityred, thick, domain=0.6:3.4, samples=40]
  plot (\x, {0.3+1.35*exp(-((\x-2.0)^2)/(2*0.42^2))});
\node[align=center] at (1.9,-0.15) {one $\xbar$ per sample, many\\ samples: $\xbar$ varies too};
\end{scope}
\begin{scope}[shift={(8.4,0)}]
\node[panel] at (1.9,1.0) {};
\node[font=\scriptsize\bfseries] at (1.9,2.5) {3. $\SE=\sigma/\sqrt{n}$};
\draw[->] (0.6,0.3) -- (3.7,0.3);
\draw[->] (0.6,0.3) -- (0.6,2.2) node[left] {SE};
\draw[distblue, thick, domain=1.13:3.6, samples=50]
  plot (\x, {0.3+1.339/sqrt(\x-0.6)});
\draw[gray, dashed] (1.3,0.3) -- (1.3,1.9);
\draw[gray, dashed] (3.4,0.3) -- (3.4,1.1);
\fill[densityred] (1.3,1.9) circle (0.05);
\fill[densityred] (3.4,1.1) circle (0.05);
\node[anchor=west] at (1.36,1.98) {SE at $n$};
\node[anchor=south] at (3.3,1.2) {SE$/2$ at $4n$};
\node[below] at (1.3,0.3) {$n$};
\node[below] at (3.4,0.3) {$4n$};
\node at (1.9,-0.3) {$4\times$ the data, half the SE};
\end{scope}
\begin{scope}[shift={(2.1,-3.75)}]
\node[panel] at (1.9,1.0) {};
\node[font=\scriptsize\bfseries] at (1.9,2.5) {4. Uncertainty intervals};
\draw[gray, dashed] (1.9,0.25) -- (1.9,2.2);
\node[gray!70!black, anchor=west] at (1.95,2.2) {truth};
\foreach \y/\c/\w in {0.45/1.75/0.6,0.75/2.05/0.55,1.05/1.6/0.65,
  1.35/2.9/0.5,1.65/2.1/0.6,1.95/1.85/0.55} {
  \pgfmathsetmacro{\lo}{\c-\w}\pgfmathsetmacro{\hi}{\c+\w}
  \ifdim \lo pt>1.9pt \def\col{trapred}\else\def\col{targetgreen}\fi
  \draw[\col, thick, |-|] (\lo,\y) -- (\hi,\y);
  \fill[\col] (\c,\y) circle (0.05);
}
\node[align=center] at (1.9,-0.15) {{\color{targetgreen}green}: covers the truth\\
  {\color{trapred}red}: misses (1 of 6)};
\end{scope}
\begin{scope}[shift={(6.3,-3.75)}]
\node[panel] at (1.9,1.0) {};
\node[font=\scriptsize\bfseries] at (1.9,2.5) {5. Signal $+$ noise};
\draw[->] (0.3,0.3) -- (3.6,0.3) node[above left] {$x$};
\draw[->] (0.3,0.3) -- (0.3,2.2) node[below right] {$y$};
\draw[structurepurple, thick] (0.5,0.55) -- (3.4,1.95);
\foreach \x/\y in {0.7/0.95,1.1/0.6,1.5/1.3,1.9/0.95,2.3/1.75,2.7/1.35,
  3.1/2.05} {
  \pgfmathsetmacro{\f}{0.55+(\x-0.5)*1.4/2.9}
  \draw[residgray] (\x,\y) -- (\x,\f);
  \fill[sourceblue] (\x,\y) circle (0.06);
}
\node[align=center] at (1.9,-0.15) {purple line: signal $f(x)$\\
  grey vertical gaps: noise $\varepsilon$};
\end{scope}
\end{tikzpicture}
\caption{The five ideas. Data vary (1); so does any statistic computed from
them, and its distribution over repeated samples is the sampling distribution
(2); its spread shrinks like $1/\sqrt{n}$ (3); that spread sets the width of
an interval that usually, but not always, covers the truth (4); and a model
splits each observation into a systematic part and leftover noise (5).}
\label{fig:pre-five-ideas}
\end{figure}

Idea 3 decides how much data every experiment needs, so here are the numbers.
Suppose single request latencies have standard deviation $\sigma = 40$ ms.
The standard error of the mean of $n$ requests is $40/\sqrt{n}$:
\[
n=25:\ 8\text{ ms},\qquad n=100:\ 4\text{ ms},\qquad
n=400:\ 2\text{ ms},\qquad n=1600:\ 1\text{ ms}.
\]
Each halving of the SE costs four times as many requests.

\begin{recoverybox}
\textbf{Pause and think.} With $\sigma = 40$ ms, how many requests do you
need for an SE of 0.5 ms? (Solve $40/\sqrt{n} = 0.5$: $\sqrt{n} = 80$, so
$n = 6400$. Four times the 1 ms answer, as the rule predicts.)
\end{recoverybox}

Figure~\ref{fig:pre-se-curve} plots the curve. A 5 ms shift in the mean is
invisible against an SE of 8 ms but stands out clearly against an SE of 1 ms.

\begin{figure}[H]
\centering
\begin{tikzpicture}
\begin{axis}[width=10cm, height=5.6cm, xmin=-40, xmax=1700, ymin=0, ymax=10,
  xlabel={number of requests $n$}, ylabel={$\SE$ of mean (ms)},
  xtick={100,400,1600}, ytick={0,1,2,4,5,8}, grid=major, grid style={gray!20},
  tick label style={font=\scriptsize}, label style={font=\small}]
\addplot[distblue, thick, domain=16:1700, samples=160] {40/sqrt(x)};
\addplot[quotientorange, thick, dashed, domain=-40:1700] {5};
\node[anchor=south east, font=\scriptsize, quotientorange!80!black]
  at (axis cs:1690,5.1) {5 ms: the change we want to detect};
\addplot[only marks, mark=*, densityred] coordinates
  {(25,8) (100,4) (400,2) (1600,1)};
\node[anchor=south west, font=\scriptsize] at (axis cs:50,8.05) {$n=25$: 8 ms};
\node[anchor=south west, font=\scriptsize] at (axis cs:110,4.1) {$n=100$: 4 ms};
\node[anchor=south west, font=\scriptsize] at (axis cs:410,2.1) {$n=400$: 2 ms};
\node[anchor=south east, font=\scriptsize] at (axis cs:1640,1.2) {$n=1600$: 1 ms};
\end{axis}
\end{tikzpicture}
\caption{Standard error of a mean latency, $40/\sqrt{n}$, when $\sigma = 40$ ms.
Where the blue curve is above the dashed line (fewer than 64 requests), the
mean's own noise is bigger than the 5 ms change we hope to see; at 400
requests the SE is 2 ms, so a 5 ms change is 2.5 SEs. The curve flattens:
going from 400 to 1600 requests buys only 1 ms more precision.}
\label{fig:pre-se-curve}
\end{figure}

\section*{Relation to \emph{Probability From the Ground Up}}

This book has a sibling, \emph{Probability From the Ground Up}. Probability
starts from a known model and asks what data it will produce. Statistics runs
the other way: it starts from data and asks what model produced it.

You do not need the probability book to read this one. What you need is a
working sense of three facts about a random quantity $X$, and each is
re-explained at the point where it is first used:
\begin{itemize}
  \item the expectation $\E[X]$ is the long-run average value of $X$;
  \item the variance $\Var(X) = \E\big[(X-\E[X])^2\big]$ measures spread, and
    $\SD(X) = \sqrt{\Var(X)}$;
  \item for constants $a,b$, $\E[aX+b] = a\E[X]+b$ and
    $\Var(aX+b) = a^2\Var(X)$; for independent $X,Y$,
    $\Var(X+Y) = \Var(X)+\Var(Y)$.
\end{itemize}
The variance rules in the last item give $\SE = \sigma/\sqrt{n}$: the sum of
$n$ independent values has variance $n\sigma^2$, and dividing by $n$ scales
that by $1/n^2$, leaving $\sigma^2/n$ (Chapter 6). When a proof
needs more machinery, such as the Central Limit Theorem or maximum likelihood
theory, we state the result, show it numerically, and point to the relevant
chapter of the probability book (Chapters 8--9 for expectation and variance,
20--21 for the law of large numbers and the CLT, 30--32 for estimation and
classical inference).

\section*{How the book is written}

Each idea starts with a small dataset worked by hand, then the general
formula, then the reason the formula is what it is. Figures carry much of the
argument; read them as carefully as the text. Coloured boxes mark specific
voices: \emph{Intuition} for the picture, \emph{Why it works} for the
derivation, \emph{Common trap} for the mistakes people actually make,
\emph{Connection} for links to machine learning, and \emph{Pause} for a quick
check to try before reading on. Short Python snippets
(\texttt{numpy}, \texttt{scipy}) show how to do each computation in practice.
Chapters 1--15 each end with a summary and exercises with answers.

Do the exercises. Statistics is easy to nod along to and hard to apply, and
the gap closes only by computing a few intervals and tests yourself.
